% メールに添付する図(研究ノート E9 の窓の図を単体にしたもの)
% ビルド: cd share && lualatex window_email.tex
% データ: ../research/figs/data/e09_window_{opt,real}.dat(calc/e09_review.py が生成)
\documentclass[border=8pt]{standalone}
\usepackage{luatexja}
\usepackage{amsmath}
\usepackage[dvipsnames]{xcolor}
\usepackage{pgfplots}
\pgfplotsset{compat=1.18}
\usepgfplotslibrary{groupplots}
\usetikzlibrary{calc}
\usepackage{varwidth}

\definecolor{cA}{HTML}{2A78D6}
\definecolor{cB}{HTML}{EB6834}
\definecolor{cC}{HTML}{1BAF7A}
\definecolor{ink}{HTML}{0B0B0B}
\definecolor{ink2}{HTML}{52514E}
\definecolor{gridc}{HTML}{E4E3DF}
\pgfplotsset{note/.style={
  axis line style={ink2}, tick style={ink2}, tick label style={font=\scriptsize, color=ink2},
  label style={font=\small, color=ink}, major grid style={gridc}, grid=major, clip=true}}
\newcommand{\Msun}{\ensuremath{\,M_\odot}}
\newcommand{\D}{../research/figs/data}

\begin{document}
\begin{varwidth}{16.5cm}
\centering
{\large 自己相互作用のある ALP で三つの観測が両立する窓}\par\medskip
\begin{tikzpicture}
\begin{groupplot}[group style={group size=2 by 1, horizontal sep=14mm},
  width=8.2cm, height=7.4cm, note, xmode=log, ymode=log,
  xmin=0.4e-12, xmax=1.2e-12, ymin=3e10, ymax=1e13,
  xtick={0.4e-12,0.6e-12,0.8e-12,1e-12}, xticklabels={0.4,0.6,0.8,1.0},
  xlabel={ALP の質量 $\mu\ [10^{-12}\,\mathrm{eV}]$}]
\nextgroupplot[ylabel={崩壊定数 $f\ [\mathrm{GeV}]$}, title={\small (a) 楽観的な仮定},
  legend to name=wleg, legend columns=2, legend style={font=\scriptsize, draw=none, column sep=14pt}, legend cell align=left]
  \fill[cB!12] (axis cs:0.60e-12,3e10) rectangle (axis cs:0.89e-12,1e13);
  \addplot[ink, line width=1.1pt] table[x=mu, y=GRS] {\D/e09_window_opt.dat}; \addlegendentry{GRS 1915+105 が回転を保つ上限}
  \addplot[cA, line width=1.1pt] table[x=mu, y=GW241011] {\D/e09_window_opt.dat}; \addlegendentry{GW241011 が回転を保つ上限}
  \addplot[cC!60!black, dashed, line width=1pt] table[x=mu, y=M15_d18] {\D/e09_window_opt.dat}; \addlegendentry{1 世代目が回転を失う下限: $15\Msun$}
  \addplot[cC, dashed, line width=1pt] table[x=mu, y=M20_d18] {\D/e09_window_opt.dat}; \addlegendentry{1 世代目が回転を失う下限: $20\Msun$}
  \addplot[cC!45, dashed, line width=1pt] table[x=mu, y=M25_d18] {\D/e09_window_opt.dat}; \addlegendentry{1 世代目が回転を失う下限: $25\Msun$}
  \node[font=\scriptsize, color=cB!80!black, anchor=north] at (axis cs:0.73e-12,9e12) {Kou+ の質量};
\nextgroupplot[title={\small (b) 現実的な仮定}]
  \fill[cB!12] (axis cs:0.60e-12,3e10) rectangle (axis cs:0.89e-12,1e13);
  \addplot[ink, line width=1.1pt] table[x=mu, y=GRS] {\D/e09_window_real.dat};
  \addplot[cA, line width=1.1pt] table[x=mu, y=GW241011] {\D/e09_window_real.dat};
  \addplot[cC!60!black, dashed, line width=1pt] table[x=mu, y=M15_d18] {\D/e09_window_real.dat};
  \addplot[cC, dashed, line width=1pt] table[x=mu, y=M20_d18] {\D/e09_window_real.dat};
  \addplot[cC!45, dashed, line width=1pt] table[x=mu, y=M25_d18] {\D/e09_window_real.dat};
  \node[font=\scriptsize, color=cB!80!black, anchor=north] at (axis cs:0.73e-12,9e12) {Kou+ の質量};
\end{groupplot}
\node[anchor=north] at ($(group c1r1.south)!0.5!(group c2r1.south)+(0,-1.0cm)$) {\pgfplotslegendfromname{wleg}};
\end{tikzpicture}

\medskip
\begin{minipage}{16cm}
\footnotesize\raggedright
実線より下で GW241011 の主星と GRS 1915+105 が観測されたスピンを保ち、破線より上で 1 世代目の BH(生まれたときのスピン 0.7)が 6 Gyr のうちに 211 の臨界スピン $+0.18$ まで回転を失う。両立する窓は、両方の実線の下で、かつ破線の上にある。橙の帯は Kou+ (arXiv:2609.02678) の質量(90\%)。

(a) GRS 1915+105 は $4.5\times10^7$ yr のあいだスピン 0.95 を保ち、GW241011 の主星は 0.78 で生まれて $3\times10^7$ yr のあいだ 0.69 を保つ。\quad
(b) GRS 1915+105 は BH の年齢 3 Gyr のあいだ 0.87 を保ち(系統誤差込みの下限 0.92 から、軽い伴星からの降着で戻せる 0.05 を差し引く)、GW241011 の主星は 0.72 で生まれて $5\times10^7$ yr のあいだ 0.69 を保つ。

二準位の準平衡、相対論的な超放射の率、自己相互作用の率は主要項(相対論的な補正を入れると、上限・下限とも 1.2--2 倍上がる)。未発表。
\end{minipage}
\end{varwidth}
\end{document}
